Integer condition codes are held in the FLAGS register. Most integer operations either update
the Z, N, C, and V bits from a temporary result, update only a subset, or leave FLAGS unchanged.
Unless an instruction description says otherwise, unmentioned FLAGS bits are preserved.

\subsection{Flag Meanings}
\begingroup\footnotesize
\begin{longtable}{@{}p{0.55in}p{1.15in}p{3.75in}@{}}
\toprule
\textbf{Bit} & \textbf{Name} & \textbf{Meaning}\\
\midrule
\endhead
\texttt{Z} & zero & Set when the result value is zero.\\
\texttt{N} & negative & Set from the most significant bit of the result at the operand size.\\
\texttt{C} & carry/borrow & Set on unsigned carry out for addition, or unsigned borrow for subtraction.\\
\texttt{V} & overflow & Set on signed overflow, or by instructions that explicitly report exceptional conditions through V.\\
\bottomrule
\end{longtable}
\manualtablecaption{Integer Condition-Code Bits}
\endgroup

\subsection{Common Integer Rules}
For an operand width of \(n\) bits, arithmetic results are evaluated modulo \(2^n\) unless the
instruction explicitly defines a wider intermediate. The sign bit is bit \(n-1\). The formulas
below describe the ordinary integer ALU rules; individual instruction pages remain authoritative
for special cases.

\begingroup\footnotesize
\begin{longtable}{@{}p{1.15in}p{4.25in}@{}}
\toprule
\textbf{Operation Class} & \textbf{FLAGS Result}\\
\midrule
\endhead
ADD, ADC & Z and N come from the stored result. C is the unsigned carry out. V is set when the source and destination signs match but the result sign differs.\\
SUB, SBB, CMP & Z and N come from the subtraction result. C is the unsigned borrow. V is set when the minuend and subtrahend signs differ and the result sign differs from the minuend. CMP does not store the temporary result. Address-destination ADD/SUB forms do not update FLAGS.\\
AND, OR, XOR, TEST & Z and N come from the logical result. C and V are cleared. TEST does not store the temporary result.\\
INC, DEC & Use the ADD or SUB rule with an implicit operand of one when the instruction form updates FLAGS. INCN and DECN leave FLAGS unchanged.\\
NEG & Computes \(0 - dst\); FLAGS follow the subtraction rule.\\
ABS & Updates FLAGS from the absolute-value result and reports overflow if the most negative signed value cannot be represented.\\
CLR & Writes zero and leaves FLAGS unchanged in the current Bedrock definition.\\
Bounds checks & Set V when the value is outside the encoded interval; Z, N, and C are unchanged.\\
\bottomrule
\end{longtable}
\manualtablecaption{Common Integer Flag Computation}
\endgroup

\subsection{Shift, Rotate, and Bit Operations}
Shift and rotate counts are read as unsigned values. Let \(W\) be the selected operand width
in bits. A zero shift count leaves the destination and FLAGS unchanged. For rotates, the
architectural count is first reduced to an effective count; an effective count of zero also
leaves the destination and FLAGS unchanged.

\begingroup\footnotesize
\begin{longtable}{@{}p{1.25in}p{4.15in}@{}}
\toprule
\textbf{Operation} & \textbf{Count Rule}\\
\midrule
\endhead
\texttt{SHL}, \texttt{SHR}, \texttt{SAR} & The count is not reduced modulo \(W\). Counts greater than \(W\) are well-defined: \texttt{SHL} and \texttt{SHR} produce zero; \texttt{SAR} produces sign fill.\\
\texttt{ROL}, \texttt{ROR} & The effective count is \(\mathit{count} \bmod W\). Effective count zero leaves the destination and FLAGS unchanged.\\
\texttt{RCL}, \texttt{RCR} & The carry bit is part of the rotate ring. The effective count is \(\mathit{count} \bmod (W + 1)\). Effective count zero leaves the destination and FLAGS unchanged.\\
\bottomrule
\end{longtable}
\manualtablecaption{Shift and Rotate Count Rules}
\endgroup

\begingroup\footnotesize
\begin{longtable}{@{}p{0.75in}p{4.65in}@{}}
\toprule
\textbf{Operation} & \textbf{FLAGS Result for Nonzero Effective Count}\\
\midrule
\endhead
\texttt{SHL} & Z and N come from the result. C is the last bit shifted out; if the count is greater than \(W\), C is zero. V is set when the signed source value multiplied by \(2^{count}\) is not representable in the selected width.\\
\texttt{SHR} & Z and N come from the result. C is the last bit shifted out; if the count is greater than \(W\), C is zero. V is cleared.\\
\texttt{SAR} & Z and N come from the result. C is the last bit shifted out; if the count is greater than \(W\), C is the original sign bit. V is cleared.\\
\texttt{ROL} & Z and N come from the result. C is the least significant result bit. V is cleared.\\
\texttt{ROR} & Z and N come from the result. C is the most significant result bit. V is cleared.\\
\texttt{RCL}, \texttt{RCR} & Z and N come from the result. C is the carry slot after the through-carry rotate. V is cleared.\\
\bottomrule
\end{longtable}
\manualtablecaption{Shift and Rotate Flag Computation}
\endgroup

Bit test/set/clear/change instructions observe the selected bit before modifying it. They set
Z when that old bit was zero and leave N, C, and V unchanged. BTEST does not modify its
destination; BSET, BCLR, and BCHG apply the requested write after recording the old-bit result
in Z.

\subsection{Conditional Consumers}
Conditional control and conditional move instructions read FLAGS but do not modify them. The
condition-code table defines how each condition mnemonic interprets Z, N, C, and V. Unconditional
aliases such as JMP and TRAP correspond to the true condition and use the unconditional spelling
as the canonical disassembly.

\subsection{Floating-Point Interactions}
Floating-point comparison and test instructions update the integer FLAGS so that ordinary
conditional branches can consume the result. An unordered floating-point comparison sets V.
Floating-point exception/status bits are tracked separately in FFLAGS and are described on the
floating-point instruction pages.
